\documentclass[12pt,a4paper]{article}

% ---- Encoding & Font ----
\usepackage{iftex}
\ifPDFTeX
  \usepackage[utf8]{inputenc}
  \usepackage[T1]{fontenc}
\else
  \usepackage{fontspec}  % XeTeX/LuaTeX (tectonic): Unicode nativo (·, —, ¿, ª, §)
\fi
\usepackage[spanish]{babel}
\usepackage{csquotes}

% ---- Page layout ----
\usepackage[top=2.5cm, bottom=2.5cm, left=2.5cm, right=2.5cm]{geometry}
\usepackage{setspace}
\onehalfspacing

% ---- Math ----
\usepackage{amsmath, amssymb, amsthm}
\usepackage{mathtools}
\usepackage{physics}
\usepackage{esint}  % \oiint

% ---- Graphics ----
\usepackage{graphicx}
\usepackage{tikz}
\usepackage{float}
\usepackage{caption}

% ---- Tables ----
\usepackage{array}
\usepackage{booktabs}
\usepackage{tabularx}
\usepackage{colortbl}

% ---- Colours ----
\usepackage{xcolor}
\definecolor{notebg}{HTML}{FFF3CD}
\definecolor{noteborder}{HTML}{FFC107}
\definecolor{boxred}{HTML}{F8D7DA}
\definecolor{boxredborder}{HTML}{F5C6CB}

% ---- Boxes ----
\usepackage{mdframed}
\usepackage{tcolorbox}
\tcbuselibrary{most}

\newtcolorbox{keybox}[1][]{
  colback=blue!5,
  colframe=blue!70!black,
  arc=4pt,
  boxrule=1pt,
  fonttitle=\bfseries,
  title={#1}
}

\newtcolorbox{warnbox}[1][]{
  colback=boxred,
  colframe=boxredborder,
  coltitle=black!75,
  arc=4pt,
  boxrule=1pt,
  fonttitle=\bfseries,
  title={#1}
}

\newtcolorbox{notebox}[1][]{
  colback=notebg,
  colframe=noteborder,
  coltitle=black!75,
  arc=4pt,
  boxrule=1pt,
  fonttitle=\bfseries,
  title={#1}
}

% ---- Hyperlinks & TOC ----
\usepackage[hidelinks]{hyperref}
\usepackage{bookmark}

% ---- Enumerate ----
\usepackage{enumitem}

% ---- Miscellaneous ----
\usepackage{icomma}
\usepackage{siunitx}
\usepackage{textcomp}
\usepackage[bottom]{footmisc}

% ---- Diagramas (circuitos, gráficas y esquemas) ----
\usepackage[europeanresistors]{circuitikz}
\usepackage{pgfplots}
\pgfplotsset{compat=1.18}
\usetikzlibrary{babel}  % babel español activa `>`; esto evita romper las flechas `->`
% courses/ElecMag2024/Clases/assets/tikzstyles.tex
% ---------------------------------------------------------------------------
% Estilos compartidos para los diagramas del curso Electromagnetismo (2024).
% Fuente única del "look" visual (colores, grosores, escalas) para que todas
% las figuras (TikZ / pgfplots / CircuiTikz) sean consistentes entre clases.
%
% Es una COPIA de courses/Fisica3-2015/Clases/assets/tikzstyles.tex: mismo
% dominio, misma paleta. Copia y no symlink para que cada curso sea
% autocontenido y la paleta de Física III pueda evolucionar aparte.
%
% Uso: la clase debe cargar antes `tikz` y `pgfplots` (y `circuitikz` si la
%      clase tiene circuitos), y luego % courses/ElecMag2024/Clases/assets/tikzstyles.tex
% ---------------------------------------------------------------------------
% Estilos compartidos para los diagramas del curso Electromagnetismo (2024).
% Fuente única del "look" visual (colores, grosores, escalas) para que todas
% las figuras (TikZ / pgfplots / CircuiTikz) sean consistentes entre clases.
%
% Es una COPIA de courses/Fisica3-2015/Clases/assets/tikzstyles.tex: mismo
% dominio, misma paleta. Copia y no symlink para que cada curso sea
% autocontenido y la paleta de Física III pueda evolucionar aparte.
%
% Uso: la clase debe cargar antes `tikz` y `pgfplots` (y `circuitikz` si la
%      clase tiene circuitos), y luego % courses/ElecMag2024/Clases/assets/tikzstyles.tex
% ---------------------------------------------------------------------------
% Estilos compartidos para los diagramas del curso Electromagnetismo (2024).
% Fuente única del "look" visual (colores, grosores, escalas) para que todas
% las figuras (TikZ / pgfplots / CircuiTikz) sean consistentes entre clases.
%
% Es una COPIA de courses/Fisica3-2015/Clases/assets/tikzstyles.tex: mismo
% dominio, misma paleta. Copia y no symlink para que cada curso sea
% autocontenido y la paleta de Física III pueda evolucionar aparte.
%
% Uso: la clase debe cargar antes `tikz` y `pgfplots` (y `circuitikz` si la
%      clase tiene circuitos), y luego \input{../assets/tikzstyles.tex}
% (compilar desde el directorio de la clase para que la ruta relativa resuelva).
% ---------------------------------------------------------------------------

% ---- Paleta (alineada con las cajas del preámbulo: keybox/notebox/warnbox) ----
\definecolor{figblue}{HTML}{1F4E79}   % azul principal (líneas, keybox)
\definecolor{figred}{HTML}{C0392B}    % rojo de énfasis (warnbox)
\definecolor{figamber}{HTML}{B8860B}  % ámbar (notebox)
\definecolor{figgray}{HTML}{555555}   % auxiliares, ejes, guías

% ---- CircuiTikz: escala y grosor de los componentes ----
% Guarda \ifdefined: la electrostática (Clases 1-14) no carga `circuitikz`, y
% sin la guarda el \input revienta con `Undefined control sequence \ctikzset`.
% Cuando el curso llegue a circuitos, el bloque se activa solo.
\ifdefined\ctikzset
  \ctikzset{
    bipoles/thickness=1,
    resistors/scale=0.9,
    capacitors/scale=0.9,
  }
\fi

% ---- TikZ: estilos reutilizables ----
% Nota: las puntas se declaran como `-{Latex[...]}` y no como `->`. La forma
% con llaves NO contiene un `>` literal, así que es inmune al `>` activo que
% introduce babel español (ver CLAUDE.md §7.3.3).
\usetikzlibrary{arrows.meta,calc,angles,quotes}

\tikzset{
  figline/.style={figblue, line width=0.9pt},
  figaux/.style={figgray, dashed, line width=0.6pt},
  % vectores
  figvec/.style={figblue, line width=0.9pt, -{Latex[length=2.2mm,width=1.6mm]}},
  figvecres/.style={figred, line width=1.3pt, -{Latex[length=2.6mm,width=1.9mm]}},
  figvecaux/.style={figgray, line width=0.7pt, -{Latex[length=1.8mm,width=1.3mm]}},
  figaxis/.style={figgray, line width=0.6pt, -{Latex[length=2mm,width=1.4mm]}},
  % cargas puntuales
  figcarga/.style={circle, inner sep=0pt, minimum size=5.4pt, draw=none},
  figpos/.style={figcarga, fill=figred},
  figneg/.style={figcarga, fill=figblue},
  figneutra/.style={figcarga, fill=figgray},
  % etiquetas
  figlbl/.style={font=\small},
  figlblaux/.style={font=\footnotesize, figgray},
}

% ---- pgfplots: estilo base de las gráficas ----
\pgfplotsset{
  emfig/.style={
    width=11cm, height=6.5cm,
    axis lines=left,
    xlabel style={font=\small},
    ylabel style={font=\small},
    tick label style={font=\footnotesize},
    every axis plot/.append style={line width=1pt},
    grid=major,
    grid style={figgray!20},
    legend cell align=left,
    legend style={font=\footnotesize, draw=figgray!40},
  },
}

% (compilar desde el directorio de la clase para que la ruta relativa resuelva).
% ---------------------------------------------------------------------------

% ---- Paleta (alineada con las cajas del preámbulo: keybox/notebox/warnbox) ----
\definecolor{figblue}{HTML}{1F4E79}   % azul principal (líneas, keybox)
\definecolor{figred}{HTML}{C0392B}    % rojo de énfasis (warnbox)
\definecolor{figamber}{HTML}{B8860B}  % ámbar (notebox)
\definecolor{figgray}{HTML}{555555}   % auxiliares, ejes, guías

% ---- CircuiTikz: escala y grosor de los componentes ----
% Guarda \ifdefined: la electrostática (Clases 1-14) no carga `circuitikz`, y
% sin la guarda el \input revienta con `Undefined control sequence \ctikzset`.
% Cuando el curso llegue a circuitos, el bloque se activa solo.
\ifdefined\ctikzset
  \ctikzset{
    bipoles/thickness=1,
    resistors/scale=0.9,
    capacitors/scale=0.9,
  }
\fi

% ---- TikZ: estilos reutilizables ----
% Nota: las puntas se declaran como `-{Latex[...]}` y no como `->`. La forma
% con llaves NO contiene un `>` literal, así que es inmune al `>` activo que
% introduce babel español (ver CLAUDE.md §7.3.3).
\usetikzlibrary{arrows.meta,calc,angles,quotes}

\tikzset{
  figline/.style={figblue, line width=0.9pt},
  figaux/.style={figgray, dashed, line width=0.6pt},
  % vectores
  figvec/.style={figblue, line width=0.9pt, -{Latex[length=2.2mm,width=1.6mm]}},
  figvecres/.style={figred, line width=1.3pt, -{Latex[length=2.6mm,width=1.9mm]}},
  figvecaux/.style={figgray, line width=0.7pt, -{Latex[length=1.8mm,width=1.3mm]}},
  figaxis/.style={figgray, line width=0.6pt, -{Latex[length=2mm,width=1.4mm]}},
  % cargas puntuales
  figcarga/.style={circle, inner sep=0pt, minimum size=5.4pt, draw=none},
  figpos/.style={figcarga, fill=figred},
  figneg/.style={figcarga, fill=figblue},
  figneutra/.style={figcarga, fill=figgray},
  % etiquetas
  figlbl/.style={font=\small},
  figlblaux/.style={font=\footnotesize, figgray},
}

% ---- pgfplots: estilo base de las gráficas ----
\pgfplotsset{
  emfig/.style={
    width=11cm, height=6.5cm,
    axis lines=left,
    xlabel style={font=\small},
    ylabel style={font=\small},
    tick label style={font=\footnotesize},
    every axis plot/.append style={line width=1pt},
    grid=major,
    grid style={figgray!20},
    legend cell align=left,
    legend style={font=\footnotesize, draw=figgray!40},
  },
}

% (compilar desde el directorio de la clase para que la ruta relativa resuelva).
% ---------------------------------------------------------------------------

% ---- Paleta (alineada con las cajas del preámbulo: keybox/notebox/warnbox) ----
\definecolor{figblue}{HTML}{1F4E79}   % azul principal (líneas, keybox)
\definecolor{figred}{HTML}{C0392B}    % rojo de énfasis (warnbox)
\definecolor{figamber}{HTML}{B8860B}  % ámbar (notebox)
\definecolor{figgray}{HTML}{555555}   % auxiliares, ejes, guías

% ---- CircuiTikz: escala y grosor de los componentes ----
% Guarda \ifdefined: la electrostática (Clases 1-14) no carga `circuitikz`, y
% sin la guarda el \input revienta con `Undefined control sequence \ctikzset`.
% Cuando el curso llegue a circuitos, el bloque se activa solo.
\ifdefined\ctikzset
  \ctikzset{
    bipoles/thickness=1,
    resistors/scale=0.9,
    capacitors/scale=0.9,
  }
\fi

% ---- TikZ: estilos reutilizables ----
% Nota: las puntas se declaran como `-{Latex[...]}` y no como `->`. La forma
% con llaves NO contiene un `>` literal, así que es inmune al `>` activo que
% introduce babel español (ver CLAUDE.md §7.3.3).
\usetikzlibrary{arrows.meta,calc,angles,quotes}

\tikzset{
  figline/.style={figblue, line width=0.9pt},
  figaux/.style={figgray, dashed, line width=0.6pt},
  % vectores
  figvec/.style={figblue, line width=0.9pt, -{Latex[length=2.2mm,width=1.6mm]}},
  figvecres/.style={figred, line width=1.3pt, -{Latex[length=2.6mm,width=1.9mm]}},
  figvecaux/.style={figgray, line width=0.7pt, -{Latex[length=1.8mm,width=1.3mm]}},
  figaxis/.style={figgray, line width=0.6pt, -{Latex[length=2mm,width=1.4mm]}},
  % cargas puntuales
  figcarga/.style={circle, inner sep=0pt, minimum size=5.4pt, draw=none},
  figpos/.style={figcarga, fill=figred},
  figneg/.style={figcarga, fill=figblue},
  figneutra/.style={figcarga, fill=figgray},
  % etiquetas
  figlbl/.style={font=\small},
  figlblaux/.style={font=\footnotesize, figgray},
}

% ---- pgfplots: estilo base de las gráficas ----
\pgfplotsset{
  emfig/.style={
    width=11cm, height=6.5cm,
    axis lines=left,
    xlabel style={font=\small},
    ylabel style={font=\small},
    tick label style={font=\footnotesize},
    every axis plot/.append style={line width=1pt},
    grid=major,
    grid style={figgray!20},
    legend cell align=left,
    legend style={font=\footnotesize, draw=figgray!40},
  },
}


% ---- Title ----
\title{\textbf{Clase 13: Coeficientes de potencial, condensadores y fuerzas} \\[0.3em]
        \large{Simetría de los $P_{ij}$, capacitancia, fuerzas a partir de la energía y primeros pasos en corriente eléctrica}}
\author{Transcripto y expandido --- Curso Electromagnetismo (OpenFING)}
\date{}

\begin{document}

\maketitle
\thispagestyle{empty}
\newpage

\tableofcontents
\newpage

\section{Repaso: conductores cargados}

La clase anterior terminó con un conjunto de $n$ conductores con cargas
$Q_1, Q_2, \dots$ rodeados de vacío o de dieléctricos \textbf{lineales}. En ese
escenario el potencial de cada conductor es una combinación lineal de las cargas
de todos, con los \textbf{coeficientes de potencial} $P_{ij}$, y la energía es una
forma cuadrática en las cargas:
\[
\phi_i = \sum_{j=1}^{n} P_{ij}\,Q_j,
\qquad
U = \frac12\sum_{i=1}^{n}\sum_{j=1}^{n} P_{ij}\,Q_i\,Q_j .
\]
El ejemplo fue el condensador de dos placas aisladas de todo lo demás; más
adelante en esta clase aparece un concepto más general de condensador.

\section{Propiedades de los coeficientes de potencial}

Son tres. El docente prueba sólo la primera, «que es la que vamos a usar más».

\subsection{Simetría}

\begin{keybox}[Los coeficientes de potencial son simétricos]
\[
\boxed{\ P_{ij} = P_{ji}\ }
\]
\end{keybox}

\textbf{Prueba.} Se dejan fijas todas las cargas salvo $Q_1$, que pasa de $Q_1$ a
$Q_1+\delta Q_1$. La variación de la energía se puede calcular de dos maneras.

\textbf{Primera: derivando la forma cuadrática.} $Q_1$ aparece en los dos
índices, así que
\[
\delta U = \frac12\sum_{j} Q_j\,P_{1j}\,\delta Q_1 + \frac12\sum_{i} Q_i\,P_{i1}\,\delta Q_1 .
\]
Como $i$ y $j$ son \textbf{índices mudos}, se renombra uno y se juntan:
\[
\delta U = \frac12\sum_i \left(P_{1i}+P_{i1}\right) Q_i\,\delta Q_1 .
\]
Tiene que dar algo simétrico: es una forma cuadrática.

\textbf{Segunda: como trabajo.} El sistema es conservativo, así que $\delta U$ es
el trabajo necesario para traer $\delta Q_1$ desde el infinito hasta el
conductor~1, que es su potencial por la carga traída:
\[
\delta U = \phi_1\,\delta Q_1 = \sum_i P_{1i}\,Q_i\,\delta Q_1 .
\]
\textbf{Comparando}, y como las $Q_i$ son arbitrarias, coeficiente a coeficiente:
\[
P_{1i} = \frac12\left(P_{1i}+P_{i1}\right)
\quad\Longrightarrow\quad
P_{1i} = P_{i1}.
\]
Lo mismo se hace variando cualquier otra carga, y de ahí el resultado general.

\begin{notebox}[No es obvio]
Por definición, $P_{ij}$ mide el potencial que una carga en el conductor $j$
genera en el $i$. Que sea igual al que una carga en $i$ genera en $j$ no sale de
la definición.
\end{notebox}

\begin{notebox}[Contexto]
Es un caso particular de la \textbf{teoría de la respuesta lineal}: a los
$P_{ij}$ se les llama a veces coeficientes de respuesta, y hay un teorema general
según el cual, en sistemas en equilibrio, las funciones de respuesta están
ligadas a cantidades naturalmente simétricas (funciones de correlación). El
docente lo menciona sin desarrollarlo.
\end{notebox}

\subsection{Las otras dos propiedades}

Se enuncian y se \textbf{admiten sin demostración}: probarlas «da un poquito de
trabajo» y no se van a usar lo suficiente.
\[
P_{ij} \ge 0,
\qquad
P_{ii} \ge P_{ij} .
\]

\begin{table}[H]
  \centering
  \begin{tabularx}{\textwidth}{lX}
    \toprule
    Propiedad & Intuición \\
    \midrule
    $P_{ij}\ge 0$ & Con el cero del potencial en el infinito, agregar una carga positiva tiende a \textbf{subir} el potencial en todos lados, no a bajarlo \\
    $P_{ii}\ge P_{ij}$ & Agregar carga a un conductor modifica más su \textbf{propio} potencial que el de los demás \\
    \bottomrule
  \end{tabularx}
\end{table}

\subsection{Ejemplo: esfera conductora y conductor puntual}

Para que no parezca todo abstracto: una \textbf{esfera conductora} de radio $a$
(el conductor 1) y un \textbf{conductor puntual} (el 2), a una distancia $d$
entre centros. Hay un problema fácil y uno difícil.

\begin{figure}[H]
  \centering
  % Clase 13 — esfera conductora (1) y conductor puntual (2).
% La transcripcion usa "R" para el radio y para la distancia; aca radio a y
% distancia d, porque el resultado solo depende de d y conviene que se vea.
\begin{tikzpicture}[scale=1]

  % esfera
  \fill[figgray!25] (0,0) circle (1.3);
  \draw[figline, line width=1pt] (0,0) circle (1.3);
  \fill[figgray] (0,0) circle (1.2pt);
  \draw[figvecaux] (0,0) -- (130:1.3);
  \node[figlblaux, fill=figgray!25, inner sep=1pt] at (130:0.72) {$a$};
  \node[figlbl] at (-0.5,-0.55) {$1$};

  % conductor puntual
  \node[figpos] (P) at (5,0) {};
  \node[figlbl] at (5,0.4) {$2$};

  % distancia entre centros
  \draw[figgray, line width=0.6pt, {Latex[length=1.8mm]}-{Latex[length=1.8mm]}] (0,-1.75) -- (5,-1.75);
  \draw[figaux] (0,-1.35) -- (0,-1.9);
  \draw[figaux] (5,-0.15) -- (5,-1.9);
  \node[figlbl, fill=white, inner sep=1.5pt] at (2.5,-1.75) {$d$};

\end{tikzpicture}

  \caption{Esfera conductora de radio $a$ y conductor puntual a distancia $d$ de
  su centro. En la transcripción radio y distancia aparecen con la misma letra;
  el resultado sólo involucra $d$.}
\end{figure}

\textbf{El fácil: el potencial que genera la esfera en el punto.} Se olvida la
carga puntual y se carga sólo la esfera con $Q$. Por Gauss, afuera de una esfera
conductora el campo es el de una carga puntual en el centro, y el potencial
también. En la posición del conductor puntual:
\[
\phi_2 = \frac{Q}{4\pi\varepsilon_0 d}
\quad\Longrightarrow\quad
P_{21} = \frac{1}{4\pi\varepsilon_0 d}.
\]

\textbf{El difícil: el potencial que genera el conductor puntual en la esfera.} Si
la esfera está descargada y el conductor puntual tiene carga $q$, la esfera se
polariza con una distribución de carga difícil de calcular: habría que resolver
Laplace en una geometría complicada. Pero por la simetría:

\begin{keybox}[Resultado en dos líneas]
\[
P_{12} = P_{21} = \frac{1}{4\pi\varepsilon_0 d}
\quad\Longrightarrow\quad
\boxed{\ \phi_{\text{esfera}} = \frac{q}{4\pi\varepsilon_0 d}\ }
\]
\end{keybox}

\begin{warnbox}[Qué es y qué no es este potencial]
Es el potencial \textbf{del conductor}, que es uno solo en toda la esfera; no es
«el potencial en el centro». Las distancias de $q$ a cada punto de la esfera son
distintas, pero el potencial es el mismo en todos. Por eso tampoco importa si la
esfera es maciza o un cascarón: la carga sólo se ubica en la superficie.
\end{warnbox}

\begin{notebox}[«Es raro el resultado, pero es correcto»]
Que la esfera genere ese potencial en el punto es obvio; al revés, no lo es. Con
los coeficientes de potencial se prueba en dos líneas, y resulta más fácil
conseguir los potenciales \textbf{sólo en los conductores} que hallar el potencial
en todo el espacio.
\end{notebox}

\section{Condensadores}

Después de la parte abstracta, una que es «como un descansito»: repaso de
Física~3, pero con alguna vuelta de tuerca (sección 4.5). El curso, dice el
docente, alterna entre cosas avanzadas con ecuaciones diferenciales y repasos.

\subsection{Una definición más general}

En Física 3 un condensador es un dispositivo para almacenar energía en un campo
eléctrico, y concretamente \textbf{dos placas conductoras} a cierta distancia,
implícitamente \textbf{aisladas del resto del mundo}: lo bastante lejos de
cualquier otra carga como para que nada afecte sus potenciales.

\begin{keybox}[Definición del curso]
Un \textbf{condensador} (o capacitor) es un dispositivo compuesto por \textbf{dos
conductores} separados tales que la \textbf{diferencia de potencial entre ellos
no depende de las cargas de otros conductores cercanos}.
\end{keybox}

Contiene a la de Física 3 como caso particular ---si no hay nada cerca, no se
puede depender de ello---, pero es más útil: en un circuito real dos condensadores
pueden estar cerca, y se quiere que cada uno conserve sus propiedades. En un
circuito, además, las cargas suelen estar en placas, no «boyando» por ahí. Hay dos
maneras típicas de conseguirlo, y una es caso particular de la otra.

\subsection{Un conductor dentro del otro}

El conductor 1 tiene una \textbf{cavidad}, y adentro está el conductor 2. Afuera
hay otros, por ejemplo el 3 y el 4.

\begin{figure}[H]
  \centering
  % Clase 13 — condensador "uno dentro del otro".
% El 2 esta en la cavidad del 1; el 3 y el 4 afuera. El contenido es que la
% region encerrada por el 1 (cavidad + conductor 2) queda equipotencial cuando
% Q2 = 0, que es la situacion en que se calculan P13 y P14.
\begin{tikzpicture}[scale=1]

  % conductor 1 con cavidad (even odd)
  \fill[figgray!25, even odd rule]
    plot[smooth cycle, tension=0.7] coordinates {(-2.2,0) (-1.6,1.6) (0.3,2.0) (2.0,1.2) (2.1,-0.9) (0.6,-1.9) (-1.4,-1.6)}
    plot[smooth cycle, tension=0.7] coordinates {(-1.3,0) (-0.9,1.0) (0.4,1.2) (1.2,0.5) (1.1,-0.6) (0.1,-1.1) (-0.8,-0.9)};
  \draw[figline, line width=1pt]
    plot[smooth cycle, tension=0.7] coordinates {(-2.2,0) (-1.6,1.6) (0.3,2.0) (2.0,1.2) (2.1,-0.9) (0.6,-1.9) (-1.4,-1.6)};
  \draw[figline, line width=1pt]
    plot[smooth cycle, tension=0.7] coordinates {(-1.3,0) (-0.9,1.0) (0.4,1.2) (1.2,0.5) (1.1,-0.6) (0.1,-1.1) (-0.8,-0.9)};
  \node[figlbl] at (1.55,-1.25) {$1$};

  % conductor 2 en la cavidad
  \fill[figgray!45, draw=figblue, line width=0.9pt]
    plot[smooth cycle, tension=0.8] coordinates {(-0.4,0.1) (0,0.55) (0.55,0.3) (0.45,-0.35) (-0.15,-0.4)};
  \node[figlbl] at (0.05,0.05) {$2$};

  % conductores 3 y 4 afuera
  \fill[figgray!25, draw=figblue, line width=0.9pt]
    plot[smooth cycle, tension=0.8] coordinates {(3.3,1.3) (4.0,1.9) (4.7,1.4) (4.1,0.8)};
  \node[figlbl] at (4.0,1.35) {$3$};
  \node[figlbl] at (5.2,1.35) {$Q_3$};
  \fill[figgray!25, draw=figblue, line width=0.9pt]
    plot[smooth cycle, tension=0.8] coordinates {(3.2,-1.2) (3.9,-0.8) (4.6,-1.3) (4.0,-1.9)};
  \node[figlbl] at (3.9,-1.3) {$4$};
  \node[figlbl] at (5.15,-1.3) {$Q_4$};

\end{tikzpicture}

  \caption{El conductor 2 está en la cavidad del 1. Cuando $Q_2 = 0$ ---que es la
  situación en que se calculan $P_{13}$ y $P_{14}$--- toda la región encerrada
  por el 1 es equipotencial, y por eso $\phi_1 - \phi_2$ no depende de $Q_3$ ni de
  $Q_4$.}
\end{figure}

\textbf{Cómo se calcula $P_{13}$:} se apagan todas las cargas salvo $Q_3$ y se
mira el potencial en 1. Pero entonces, en particular, \textbf{$Q_2 = 0$}. Sin
carga adentro de la cavidad, toda la región encerrada por el conductor 1
---cavidad y conductor 2 incluidos--- tiene campo nulo y \textbf{el mismo
potencial}. Entonces $\phi_1 = \phi_2$ en esa situación, y
\[
P_{13} = P_{23},\qquad P_{14} = P_{24}.
\]
La diferencia de potencial entre 1 y 2 es
\[
\phi_1-\phi_2 = (P_{11}-P_{21})\,Q_1 + (P_{12}-P_{22})\,Q_2
+ \underbrace{(P_{13}-P_{23})}_{0}\,Q_3 + \underbrace{(P_{14}-P_{24})}_{0}\,Q_4 ,
\]
que \textbf{no depende de $Q_3$ ni de $Q_4$}. Con 28 conductores afuera sería lo
mismo.

\begin{notebox}[Receta de ingeniería]
Para un condensador «fiable», poner una placa dentro de la otra. Los
condensadores de ferretería son conductores \textbf{enrollados}: no están
totalmente uno dentro del otro, pero casi.
\end{notebox}

\subsection{Placas muy grandes}

El otro caso típico, «el caso de escuela», son \textbf{dos placas muy grandes}
comparadas con la distancia entre ellas. Es un caso particular del anterior: como
no se llegan a ver los bordes, podría ser que allá lejos las placas dieran la
vuelta y una quedara encerrada en la otra, y el resultado no puede depender del
borde.

\begin{notebox}[La analogía del docente]
Parados acá no sabemos si la Tierra es plana o esférica, porque el borde está tan
lejos que puede haberse curvado.
\end{notebox}

Las dos maneras de \textbf{apantallar} lo que pasa afuera son, entonces, una
dentro de la otra o placas muy grandes. En Física 3 este problema «se ponía abajo
de la alfombra».

\subsection{Capacitancia y energía}

Los condensadores se usan con \textbf{carga neta nula}: $Q_1 = Q$, $Q_2 = -Q$.
Usando $P_{21} = P_{12}$:
\[
\Delta\phi = \phi_1-\phi_2 = \left(P_{11} - 2P_{12} + P_{22}\right) Q .
\]
La diferencia de potencial es proporcional a la carga de una de las placas, como
se sabía de Física 3, y el coeficiente define la \textbf{capacitancia}:

\begin{keybox}[Capacitancia]
\[
\boxed{\ C = \frac{1}{P_{11}+P_{22}-2P_{12}}\ },
\qquad
\boxed{\ \Delta\phi = \frac{Q}{C}\ } .
\]
No depende de la carga sino de la \textbf{geometría} y de los
\textbf{dieléctricos} del entorno. La única hipótesis usada es que esos
dieléctricos sean \textbf{lineales}.
\end{keybox}

\textbf{Unidades.} $C = Q/\Delta\phi$ se mide en coulomb por volt, el
\textbf{farad} (F). Es una unidad enorme ---el volt es de laboratorio, el coulomb
no---: los condensadores reales son de micro-, nano- o picofarads.

\textbf{Energía.} Con la simetría, la fórmula de la clase pasada se simplifica:
\[
U = \frac12\left(P_{11}+P_{22}-P_{12}-P_{21}\right)Q^2
= \frac12\left(P_{11}+P_{22}-2P_{12}\right)Q^2
= \boxed{\ \frac{Q^2}{2C}\ } .
\]

\subsection{Condensador de placas paralelas}

Placas de área $A$ a distancia $d$, con vacío entre ellas. Por invariancia de
traslación a lo largo del plano, las densidades son uniformes, $+\sigma$ y
$-\sigma$. Por Gauss el campo entre placas es uniforme:
\[
E = \frac{\sigma}{\varepsilon_0} = \frac{Q}{A\,\varepsilon_0},
\qquad
\Delta\phi = E\,d = \frac{Q\,d}{\varepsilon_0 A}
\quad\Longrightarrow\quad
\boxed{\ C = \frac{\varepsilon_0 A}{d}\ } .
\]

\begin{notebox}
Con Laplace o el método de las imágenes ahora se pueden hallar capacitancias de
geometrías más difíciles que las de Física 3; eso queda para el práctico.
\end{notebox}

\subsection{Serie y paralelo}

Más repaso, «porque estaban oxidados».

\begin{figure}[H]
  \centering
  % Clase 13 — condensadores en serie y en paralelo (primer circuitikz del curso).
% Mismos C1 y C2 en los dos paneles; cambia solo la conexion. En (a) el
% recuadro punteado es el argumento de la clase: ese tramo esta aislado y es
% neutro, asi que las dos placas que encierra llevan -Q y +Q. Adentro solo
% rotulos; lo que se comparte y lo que se suma va en el caption.
\begin{tikzpicture}[scale=1]

  % =============== (a) serie ===============
  \begin{scope}[line width=0.9pt, color=figblue]
    \draw (0,0) to[battery1] (0,2.2);
    \draw (0,2.2) to[C] (2.4,2.2) to[C] (4.8,2.2);
    \draw (4.8,2.2) -- (4.8,0) -- (0,0);
  \end{scope}
  \node[figlbl, anchor=east] at (-0.45,1.1) {$\Delta\phi$};
  \node[figlbl, figblue] at (1.2,1.3) {$C_1$};
  \node[figlbl, figblue] at (3.6,1.3) {$C_2$};
  % cargas en las placas (placas en x = 1,1 | 1,3 y 3,5 | 3,7)
  \node[figlblaux, figred, anchor=east] at (1.1,2.95) {$+Q$};
  \node[figlblaux, figblue, anchor=west] at (1.3,2.95) {$-Q$};
  \node[figlblaux, figred, anchor=east] at (3.5,2.95) {$+Q$};
  \node[figlblaux, figblue, anchor=west] at (3.7,2.95) {$-Q$};
  % tramo aislado y neutro: encierra las dos placas interiores
  \draw[figaux] (1.2,1.65) rectangle (3.6,2.75);
  \node[figlbl, anchor=north] at (2.4,-0.35) {(a) serie};

  % =============== (b) paralelo ===============
  \begin{scope}[xshift=6.6cm]
    \begin{scope}[line width=0.9pt, color=figblue]
      \draw (0,0) to[battery1] (0,2.2) -- (4.2,2.2);
      \draw (1.8,2.2) to[C] (1.8,0);
      \draw (4.2,2.2) to[C] (4.2,0);
      \draw (4.2,0) -- (0,0);
      \fill (1.8,2.2) circle (2pt);
      \fill (1.8,0) circle (2pt);
    \end{scope}
    \node[figlbl, anchor=east] at (-0.45,1.1) {$\Delta\phi$};
    \node[figlbl, figblue, anchor=east] at (1.35,1.1) {$C_1$};
    \node[figlbl, figblue, anchor=east] at (3.75,1.1) {$C_2$};
    \node[figlblaux, figred, anchor=west] at (2.2,1.45) {$+Q_1$};
    \node[figlblaux, figblue, anchor=west] at (2.2,0.75) {$-Q_1$};
    \node[figlblaux, figred, anchor=west] at (4.6,1.45) {$+Q_2$};
    \node[figlblaux, figblue, anchor=west] at (4.6,0.75) {$-Q_2$};
    \node[figlbl, anchor=north] at (2.1,-0.35) {(b) paralelo};
  \end{scope}

\end{tikzpicture}

  \caption{Los mismos dos condensadores, conectados de dos maneras. En serie, el
  tramo punteado está aislado del resto y es neutro, así que las placas que
  encierra llevan $-Q$ y $+Q$: la carga es la misma en los dos y las diferencias
  de potencial se suman. En paralelo, la diferencia de potencial es la misma y
  las cargas se suman.}
\end{figure}

\textbf{En serie.} Dos condensadores $C_1$ y $C_2$ uno detrás del otro, con
$\Delta\phi$ aplicada a los extremos. Aparece $Q$ en una placa del primero y $-Q$
en la otra; el tramo que une los dos condensadores es neutro, así que el segundo
también tiene $\pm Q$. La diferencia de potencial total es la suma:
\[
\Delta\phi = \frac{Q}{C_1} + \frac{Q}{C_2}
\quad\Longrightarrow\quad
\boxed{\ \frac{1}{C} = \frac{1}{C_1} + \frac{1}{C_2}\ } .
\]

\textbf{En paralelo.} Unidos por cables ideales ---o simplemente en equilibrio---,
los dos tienen la \textbf{misma} $\Delta\phi$, y la carga total es la suma,
$Q = Q_1 + Q_2 = C_1\Delta\phi + C_2\Delta\phi$:
\[
\boxed{\ C = C_1 + C_2\ } .
\]

\begin{notebox}
Para quien se acuerda a medias: con las resistencias la relación es \textbf{al
revés}. Las resistencias aparecen en la clase siguiente.
\end{notebox}

\section{Fuerzas y momentos electrostáticos}

Para cerrar el capítulo 4 (sección 4.6): cómo reconstruir fuerzas y torques a
partir de la energía.

\subsection{Sistema aislado}

Un sistema electrostático \textbf{aislado}: nada le aporta energía desde afuera.
Se imagina que todo está quieto salvo una parte, que se desplaza $d\vec r$. El
trabajo electrostático es $dW = \vec F\cdot d\vec r$. Como el sistema es estático
no cambia la energía cinética, no hay fuerzas externas, y el trabajo es a costa de
la energía potencial: $dW = -dU$. Entonces

\begin{keybox}[Sistema aislado]
\[
\boxed{\ \vec F = -\grad U\ },
\qquad
\boxed{\ \tau_n = -\pdv{U}{\theta}\ } .
\]
La segunda, para una \textbf{rotación} de ángulo $d\theta$ alrededor de un eje de
versor $\hat n$, con $dW = \tau_n\,d\theta$.
\end{keybox}

Conocer la energía permite calcular no sólo fuerzas sino también torques.

\subsection{Baterías que mantienen los potenciales}

Si el sistema no está aislado «puede pasar de todo»: hay que decir cuáles son las
fuerzas externas, y el caso general no se puede tratar. Se estudia un caso
particular: \textbf{baterías} que mantienen \textbf{constantes} los potenciales de
los conductores (fijado el cero en el infinito, tiene sentido hablar de
potenciales y no sólo de diferencias).

Ahora el trabajo tiene dos partes: la asociada a la energía potencial, y la que
aporta la batería. Con $U = \tfrac12\sum_i\phi_i Q_i$ y los potenciales fijos,
sólo cambian las cargas:
\[
dU\big|_{\phi} = \frac12\sum_i \phi_i\,dQ_i .
\]
La batería mantiene los potenciales moviendo cargas $dQ_i$ a cada conductor, y el
trabajo que hace es
\[
dW_{\text{bat}} = \sum_i \phi_i\,dQ_i = 2\,dU .
\]
Esa «coincidencia grata» decide todo: el trabajo total es
$\vec F\cdot d\vec r = -dU + dW_{\text{bat}} = +dU$. Las fórmulas son las del
sistema aislado \textbf{con el signo al revés}:

\begin{keybox}[Potenciales fijos por baterías]
\[
\boxed{\ \vec F = +\grad U\big|_{\phi}\ },
\qquad
\boxed{\ \tau_n = +\pdv{U}{\theta}\bigg|_{\phi}\ } .
\]
\end{keybox}

\begin{warnbox}[No viola nada]
El sistema no está aislado: la batería aporta energía, y justo el doble de la
variación de la energía potencial. Una contribución va para un lado y la otra
para el otro.
\end{warnbox}

\begin{notebox}
El docente admite que la primera vez le resultó «re extraño», y que no conoce una
explicación más simple que la cuenta misma: a potencial fijo, el trabajo de la
batería es el doble de la variación de la energía potencial.
\end{notebox}

\subsection{Ejemplo: un dieléctrico entre las placas}

Un condensador de placas paralelas rectangulares, de largo $L$ y ancho $W$,
separadas $D$, conectado a una batería que mantiene $\Delta\phi$ fijo. Entre las
placas hay un bloque de dieléctrico (lineal, homogéneo, isótropo, constante $K$)
que puede entrar o salir: un tramo de largo $x$ queda en vacío y el resto,
$L-x$, con dieléctrico. \textbf{¿La fuerza lo chupa hacia adentro o lo expulsa?}

\begin{figure}[H]
  \centering
  % Clase 13 — dielectrico que entra en un condensador conectado a una bateria.
% Vista de costado: largo L, separacion D (el ancho W va hacia adentro de la
% hoja). Tramo x en vacio a la izquierda, L-x con dielectrico. La fuerza va hacia
% la izquierda: achica x, es decir, chupa el dielectrico.
\begin{tikzpicture}[scale=1]

  % dielectrico (sobresale a la derecha)
  \fill[figamber!20] (2.4,0.06) rectangle (7.6,1.44);
  \draw[figamber, line width=0.8pt] (2.4,0.06) rectangle (7.6,1.44);
  \node[figlbl, figamber!80!black] at (4.2,0.75) {$K$};

  % placas
  \draw[figline, line width=2pt] (0,1.5) -- (6,1.5);
  \draw[figline, line width=2pt] (0,0) -- (6,0);

  % bateria
  \draw[figline] (0,1.5) -- (-1.2,1.5) -- (-1.2,0.95);
  \draw[figline] (-1.2,0.55) -- (-1.2,0) -- (0,0);
  \draw[figline, line width=1.2pt] (-1.55,0.95) -- (-0.85,0.95);
  \draw[figline, line width=1.2pt] (-1.38,0.55) -- (-1.02,0.55);
  \node[figlbl] at (-2.05,0.75) {$\Delta\phi$};

  % fuerza sobre el dielectrico
  \draw[figvecres] (7.3,0.75) -- (6.3,0.75);
  \node[figlbl, figred] at (6.8,1.08) {$F$};

  % cotas
  \draw[figgray, line width=0.6pt, {Latex[length=1.6mm]}-{Latex[length=1.6mm]}] (0,-0.45) -- (2.4,-0.45);
  \node[figlbl, fill=white, inner sep=1pt] at (1.2,-0.45) {$x$};
  \draw[figgray, line width=0.6pt, {Latex[length=1.6mm]}-{Latex[length=1.6mm]}] (2.4,-0.45) -- (6,-0.45);
  \node[figlbl, fill=white, inner sep=1pt] at (4.2,-0.45) {$L-x$};
  \draw[figgray, line width=0.6pt, {Latex[length=1.6mm]}-{Latex[length=1.6mm]}] (0,-1.0) -- (6,-1.0);
  \node[figlbl, fill=white, inner sep=1pt] at (3,-1.0) {$L$};
  \draw[figgray, line width=0.6pt, {Latex[length=1.6mm]}-{Latex[length=1.6mm]}] (1.1,0.12) -- (1.1,1.38);
  \node[figlbl] at (1.4,0.75) {$D$};

\end{tikzpicture}

  \caption{Vista de costado; el ancho $W$ de las placas va hacia adentro de la
  hoja. Con la batería manteniendo $\Delta\phi$, la fuerza sobre el dieléctrico
  apunta hacia adentro del condensador.}
\end{figure}

Se \textbf{desprecian los efectos de borde}: si $L\gg D$, las zonas de borde son
del ancho de $D$ y el grueso de la energía está en las regiones interiores, que
es proporcional a $x$ o a $L-x$. Así el sistema son \textbf{dos condensadores en
paralelo}:
\[
C_0 = \frac{\varepsilon_0\, x\,W}{D},
\qquad
C_K = \frac{K\varepsilon_0\,(L-x)\,W}{D}.
\]
Usando $\Delta\phi = Q/C$, cada energía se escribe $\tfrac12 C\,\Delta\phi^2$:
\[
U = \frac12\left(C_0 + C_K\right)\Delta\phi^2
= \frac{\Delta\phi^2}{2}\,\frac{\varepsilon_0 W}{D}\left[x + K(L-x)\right].
\]
Con $\Delta\phi$ fijo por la batería, la fuerza es \textbf{más} la derivada:
\[
F_x = +\pdv{U}{x} = \frac{\Delta\phi^2}{2}\,\frac{\varepsilon_0 W}{D}\,(1-K) .
\]

\begin{keybox}[El condensador chupa el dieléctrico]
Como $K>1$, $F_x<0$: la fuerza tiende a achicar el tramo en vacío. El dieléctrico
es atraído hacia el interior.
\end{keybox}

\begin{notebox}[Ejercicio]
¿Qué pasa si el condensador está cargado pero \textbf{desconectado} de la
batería? Entonces $\Delta\phi$ ya no es fijo, pero la carga sí, porque nadie
aporta carga. ¿El dieléctrico tiende a entrar o a salir?
\end{notebox}

\begin{warnbox}[Por qué acá el signo es $+$]
Porque el sistema no está aislado. Si lo estuviera, sería $-\grad U$. Los dos
casos fáciles de tratar son el aislado y el de potenciales fijos.
\end{warnbox}

\section{Corriente eléctrica}

Con los minutos que quedan arranca el \textbf{capítulo 5}: por primera vez en el
curso las cargas se mueven.

\subsection{Definición}

Se toma una superficie $S$ y cargas que la cruzan. La \textbf{carga que cruza $S$
por unidad de tiempo} es la \textbf{corriente eléctrica}:
\[
\boxed{\ I = \frac{dQ}{dt}\ } .
\]

\begin{warnbox}[Convención de signos]
Se cuenta \textbf{carga}, no partículas, y hay que elegir qué sentido de cruce es
positivo. La corriente va siempre en el sentido de hipotéticas cargas
\textbf{positivas} en movimiento: si lo que se mueve son cargas negativas, la
corriente va en sentido opuesto a su movimiento.
\end{warnbox}

\textbf{Unidad:} coulomb por segundo, el \textbf{ampere} (A).

\subsection{Tipos de conducción}

\begin{enumerate}
  \item \textbf{Electrones en un metal.} Se mueven los electrones de valencia en
    una red cristalina iónica. En los semiconductores a veces se mueven
    \textbf{huecos}, ausencias de electrones que van para el otro lado; su
    naturaleza no tiene sentido clásico. Cualitativamente sirve pensar el metal
    como una caja con electrones que se mueven, aunque en realidad la red vibra y
    tiene impurezas.
  \item \textbf{Electrolitos en un fluido.} Iones disueltos en un líquido.
  \item \textbf{Transporte hidrodinámico en un plasma.} En el Sol, por ejemplo, la
    materia está tan caliente que electrones e iones están separados: es un fluido
    cargado, y lo que se mueve es \textbf{el propio medio}, no cargas dentro de
    él.
\end{enumerate}

\begin{notebox}[Lo que no se estudia]
El tercero: habría que acoplar las ecuaciones del electromagnetismo con las de la
hidrodinámica (las de Navier--Stokes para un fluido neutro), y cada una complica a
la otra. Es la única vez que el curso menciona los plasmas.
\end{notebox}

\medskip
\noindent\emph{La clase siguiente continúa con la corriente eléctrica y las
resistencias.}

\end{document}
